\documentclass[10pt,a4paper]{article}
\usepackage[margin=1.8cm]{geometry}
\usepackage{amsmath, amssymb, graphicx}
\usepackage[hidelinks]{hyperref}
\setlength{\parindent}{0pt}
\setlength{\parskip}{0.4em}
\pagestyle{empty}
% The pooled segment cloglog GLMM, explained on made-up data. Figures: cloglog_explained_figures.py.
% Text left, figure right, as in fpca_from_plots.tex.
\newcommand{\glmmrow}[2]{\par\noindent
  \begin{minipage}[t]{0.52\textwidth}\vspace{0pt}#1\end{minipage}\hfill
  \begin{minipage}[t]{0.45\textwidth}\vspace{0pt}#2\end{minipage}\par\medskip}
\newcommand{\glmmfig}[2]{\centering\includegraphics[width=\linewidth]{week_figure/cloglog/#1}\\[-0.1em]
  {\footnotesize #2}}

\begin{document}
\sloppy

\subsection*{The pooled segment cloglog GLMM}
Both the geometry and the outcome are recorded per segment: curvature and taper are measured in each
of the nine trunk segments, and each segment is read for incident coronary plaque at follow-up. The
risk model is therefore fitted on the segments, one row per segment. The who question is about the
person, so the segment risks are then pooled into one person risk. The model is a cloglog generalized linear mixed model (GLMM): a cloglog
link, a linear predictor, and a person random intercept. A person-level model of any plaque is kept as
the benchmark. The figures use made-up numbers.

\glmmrow{%
\textbf{The cloglog link.} Each segment is plaque-free at baseline and seen again about ten years
later, and when plaque started is unknown. This is grouped survival data, for which the proportional
hazards model is a cloglog model (Prentice and Gloeckler 1978). The name lists the steps: the
\emph{complement} $1-p$ is the probability that the segment stays plaque-free, the first \emph{log}
gives the hazard built up over follow-up, $\Lambda = -\log(1-p)$, and the second \emph{log} gives
$\eta = \log\Lambda$, which can take any value and is modelled as a sum of terms:
\[
  \begin{split}
  \mathrm{cloglog}\,P(y_{ij}=1 \mid b_i) = {} & \alpha_j + \beta_\mathrm{S}\,\mathrm{SCORE2}_i \\
    & + \beta_\mathrm{L}\log\ell_{ij} + \text{geometry}_{ij} + b_i ,
  \end{split}
\]
with $\ell_{ij}$ the segment length. $e^{\beta}$ is then a hazard ratio, which multiplies the hazard by the same factor in the pLAD and the
rarer dLCx. An odds ratio would drift with the segment's baseline risk.
}{%
\glmmfig{c3_hr_vs_or.png}{A hazard ratio of 1.5 is the same at every baseline risk; the odds ratio is
not.}
}

\glmmrow{%
\textbf{Hazards add, the person effect multiplies.} A person is free of plaque only if all nine
segments are, so the probabilities $e^{-\Lambda_j}$ multiply and the hazards add. The person risk
depends on the segments only through the summed hazard $S_i = \sum_j e^{\eta_{ij}}$, and an
unevaluable segment drops out of the sum. Risk that nobody measured acts on all segments of a person,
so the nine rows of one person are not independent. The random intercept $b_i \sim N(0, \sigma^2)$
captures this: it multiplies all nine hazards of person $i$ by $e^{b_i}$, and only $\sigma$ is
estimated.
}{%
\glmmfig{r2_hazards_add.png}{The nine segment hazards of one made-up person, stacked, for three values
of the person effect $b$.}
}

\glmmrow{%
\textbf{One person risk.} For a new person $b$ is unknown, so the person risk averages over it:
\[
  P_i = 1 - \int \exp\!\big(-e^{b} S_i\big)\,\phi(b; 0, \sigma^2)\,\mathrm{d}b ,
\]
computed by Gauss--Hermite quadrature. For the made-up person $P = 0.548$. Setting $b = 0$ gives 0.534,
the risk of the median person, and multiplying the average segment risks, which ignores the shared
$b$, gives 0.687. Both shortcuts rank people correctly, so only calibration detects them.
}{%
\glmmfig{r3_average_over_b.png}{Person risk for each value of $b$ (blue), the spread of $b$ (grey) and
the average used (green).}
}

\textbf{Fitting and testing.} The weights and $\sigma$ are fitted by maximum likelihood with $b$
integrated out. Cross-validation is split by person, and $P_i$ uses the $\sigma$ of the training folds.
Geometry is tested once, by a likelihood-ratio test of the geometry block; the gain in C and the
calibration are estimated, not tested. The current code (\texttt{progression\_who\_pooled.py}) runs on
the Medis trees with progression as the outcome and two geometry terms (mean turning and the LM angle).

\end{document}
